\documentclass[11pt,a4paper]{article}
\usepackage[utf8]{inputenc}
\usepackage[T1]{fontenc}
\usepackage[spanish,es-noshorthands]{babel}
\usepackage[margin=2.5cm]{geometry}
\usepackage{amsmath,amssymb,mathtools}
\usepackage{enumitem}
\usepackage{booktabs}
\usepackage{tikz}
\usepackage{pgfplots}
\pgfplotsset{compat=1.18}

\newcounter{ejnum}
\newenvironment{ejercicio}{\refstepcounter{ejnum}\par\medskip\noindent\textbf{Ejercicio \theejnum.}~}{\par}
\newenvironment{solucion}{\par\noindent\textit{Solución.}~}{\par\vspace{1em}}

\title{Práctica 2: Funciones de varias variables y distribuciones del muestreo \\ \large Unidad 3: Vectores aleatorios}
\author{Probabilidad y Estadística (5765) \\ Universidad Nacional del Sur}
\date{}

\begin{document}
\maketitle

\begin{ejercicio}
Sean $X$ e $Y$ variables aleatorias discretas independientes, con
\[
P(X=0)=P(X=1)=P(X=2)=\tfrac13, \qquad P(Y=0)=P(Y=1)=\tfrac12.
\]
Hallar la función de probabilidad de $Z=X+Y$ mediante convolución.
\end{ejercicio}

\begin{solucion}
$Z$ puede tomar los valores $0,1,2,3$. Usando $p_Z(z) = \sum_k p_X(k)\,p_Y(z-k)$:
\[
p_Z(0) = p_X(0)p_Y(0) = \tfrac13\cdot\tfrac12 = \tfrac16.
\]
\[
p_Z(1) = p_X(0)p_Y(1)+p_X(1)p_Y(0) = \tfrac13\cdot\tfrac12+\tfrac13\cdot\tfrac12 = \tfrac13.
\]
\[
p_Z(2) = p_X(1)p_Y(1)+p_X(2)p_Y(0) = \tfrac13\cdot\tfrac12+\tfrac13\cdot\tfrac12 = \tfrac13.
\]
\[
p_Z(3) = p_X(2)p_Y(1) = \tfrac13\cdot\tfrac12 = \tfrac16.
\]
Verificación: $\tfrac16+\tfrac13+\tfrac13+\tfrac16 = 1$. \checkmark
\end{solucion}

\begin{ejercicio}
Sean $X,Y\sim U(0,1)$ independientes. Hallar la densidad de $Z=X-Y$ y graficar mentalmente su forma.
\end{ejercicio}

\begin{solucion}
Usamos convolución con $f_Y(-y)$: definamos $W=-Y$, que tiene densidad $f_W(w) = f_Y(-w) = 1$ para $w\in(-1,0)$. Entonces $Z=X+W$, y por convolución,
\[
f_Z(z) = \int f_X(x)\, f_W(z-x)\, dx,
\]
con $f_X(x)=1$ en $(0,1)$ y $f_W(z-x)=1$ cuando $z-x\in(-1,0) \iff x\in(z,z+1)$. Intersecando con $x\in(0,1)$:
\begin{itemize}
\item Si $-1\le z\le 0$: $x\in(0,z+1)$, longitud $z+1$, entonces $f_Z(z)=z+1$.
\item Si $0< z\le 1$: $x\in(z,1)$, longitud $1-z$, entonces $f_Z(z)=1-z$.
\item Fuera de $[-1,1]$: $f_Z(z)=0$.
\end{itemize}
Resulta una distribución \textbf{triangular simétrica} en $[-1,1]$, con pico en $z=0$ y valor máximo $f_Z(0)=1$.
\end{solucion}

\begin{ejercicio}
El tiempo de espera $X$ en la caja 1 de un supermercado y el tiempo de espera $Y$ en la caja 2 son independientes, ambos con distribución $\text{Exp}(\lambda=0.5)$ (minutos$^{-1}$). Un cliente elige al azar una de las dos cajas... en realidad, nos interesa el tiempo total de dos clientes distintos, uno en cada caja: $Z=X+Y$.
\begin{enumerate}[label=(\alph*)]
\item Hallar la densidad de $Z$ mediante convolución.
\item Calcular $P(Z>6)$.
\end{enumerate}
\end{ejercicio}

\begin{solucion}
\textbf{(a)} Con $f_X(x)=f_Y(x)=0.5\,e^{-0.5x}$, $x>0$, para $z>0$:
\[
f_Z(z) = \int_0^z 0.5e^{-0.5x}\cdot 0.5e^{-0.5(z-x)}\,dx = 0.25\,e^{-0.5z}\int_0^z dx = 0.25\,z\,e^{-0.5z}, \quad z>0,
\]
que es una densidad $\text{Gamma}(2, 0.5)$.

\textbf{(b)} Integramos por partes (o usamos la función de distribución Gamma con $n=2$: $P(Z>z)=e^{-\lambda z}(1+\lambda z)$ para Gamma$(2,\lambda)$):
\[
P(Z>6) = e^{-0.5\cdot 6}\left(1+0.5\cdot 6\right) = e^{-3}(1+3) = 4e^{-3} \approx 0.1992.
\]
\end{solucion}

\begin{ejercicio}
Un sistema tiene 4 componentes idénticos e independientes conectados en \textbf{paralelo} (el sistema funciona mientras al menos un componente funcione). Cada componente tiene vida útil $\text{Exp}(\lambda=0.1)$ (fallas por año). Sea $M$ la vida del sistema, $M=\max(X_1,X_2,X_3,X_4)$.
\begin{enumerate}[label=(\alph*)]
\item Hallar $F_M(t)$.
\item Calcular $P(M>15)$.
\end{enumerate}
\end{ejercicio}

\begin{solucion}
\textbf{(a)} Cada $X_i$ tiene $F(t) = 1-e^{-0.1t}$, $t>0$. Por independencia,
\[
F_M(t) = \prod_{i=1}^4 F(t) = \left(1-e^{-0.1t}\right)^4, \quad t>0.
\]

\textbf{(b)}
\[
P(M>15) = 1-F_M(15) = 1-\left(1-e^{-1.5}\right)^4.
\]
Con $e^{-1.5}\approx 0.2231$: $1-0.2231 = 0.7769$, y $0.7769^4 \approx 0.3644$. Entonces $P(M>15) \approx 1-0.3644 = 0.6356$.
\end{solucion}

\begin{ejercicio}
Retomando el sistema del Ejercicio 4, pero ahora conectado en \textbf{serie} (el sistema falla apenas falla el primer componente). Sea $m^*=\min(X_1,\ldots,X_4)$.
\begin{enumerate}[label=(\alph*)]
\item Mostrar que $m^*\sim\text{Exp}(0.4)$.
\item Calcular $P(m^*>5)$ y comparar con la probabilidad análoga para un único componente, $P(X_1>5)$.
\end{enumerate}
\end{ejercicio}

\begin{solucion}
\textbf{(a)} $1-F_{m^*}(t) = \prod_{i=1}^4 (1-F(t)) = \left(e^{-0.1t}\right)^4 = e^{-0.4t}$, luego $F_{m^*}(t) = 1-e^{-0.4t}$, es decir $m^*\sim\text{Exp}(0.4)$ (suma de las tasas, como se vio en la teoría).

\textbf{(b)}
\[
P(m^*>5) = e^{-0.4\cdot5} = e^{-2} \approx 0.1353.
\]
\[
P(X_1>5) = e^{-0.1\cdot5} = e^{-0.5} \approx 0.6065.
\]
Como era de esperar, un sistema en serie con más componentes es mucho menos confiable ($0.1353 \ll 0.6065$): basta que falle uno solo de los cuatro para que el sistema completo falle.
\end{solucion}

\begin{ejercicio}
Sean $X\sim N(50,16)$ e $Y\sim N(30,9)$ independientes (por ejemplo, $X$ e $Y$ son los puntajes de dos secciones distintas de un examen). Sea $Z=X+Y$ el puntaje total.
\begin{enumerate}[label=(\alph*)]
\item Indicar la distribución de $Z$ usando la propiedad reproductiva de la normal.
\item Calcular $P(Z>90)$.
\end{enumerate}
\end{ejercicio}

\begin{solucion}
\textbf{(a)} Por la propiedad reproductiva, $Z = X+Y \sim N(50+30,\, 16+9) = N(80, 25)$, es decir, $\sigma_Z = 5$.

\textbf{(b)} Estandarizando: $P(Z>90) = P\left(\dfrac{Z-80}{5} > \dfrac{90-80}{5}\right) = P(W>2)$ con $W\sim N(0,1)$.
\[
P(W>2) = 1-\Phi(2) \approx 1-0.9772 = 0.0228.
\]
\end{solucion}

\begin{ejercicio}
Sean $Z_1,Z_2,Z_3,Z_4,Z_5 \sim N(0,1)$ independientes. Se define $U = Z_1^2+Z_2^2+Z_3^2$ y $V=Z_4^2+Z_5^2$.
\begin{enumerate}[label=(\alph*)]
\item Indicar las distribuciones de $U$, de $V$ y de $U+V$.
\item Calcular $E[U+V]$ y $\text{Var}(U+V)$.
\item Si $T = \dfrac{Z_1}{\sqrt{V/2}}$, indicar la distribución de $T$ y de $T^2$.
\end{enumerate}
\end{ejercicio}

\begin{solucion}
\textbf{(a)} Por definición, $U\sim\chi_3^2$ (suma de 3 cuadrados de normales estándar independientes) y $V\sim\chi_2^2$. Como $U$ y $V$ dependen de conjuntos disjuntos de $Z_i$'s independientes, son independientes entre sí, y por la propiedad reproductiva, $U+V\sim\chi_5^2$.

\textbf{(b)} $E[U+V] = 5$ (grados de libertad), $\text{Var}(U+V) = 2\times 5 = 10$.

\textbf{(c)} $Z_1\sim N(0,1)$ es independiente de $V\sim\chi_2^2$ (pues $Z_1$ no interviene en $V$). Por definición de la $t$ de Student, $T = Z_1/\sqrt{V/2} \sim t_2$. Además, por la relación entre $t$ y $F$, $T^2 \sim F_{1,2}$.
\end{solucion}

\begin{ejercicio}
Se toman dos muestras aleatorias independientes de tamaños $n_1=8$ y $n_2=12$ de dos poblaciones normales con varianzas $\sigma_1^2=\sigma_2^2=\sigma^2$ (desconocida pero igual). Sean $S_1^2$ y $S_2^2$ las varianzas muestrales.
\begin{enumerate}[label=(\alph*)]
\item Indicar la distribución de $\dfrac{S_1^2}{S_2^2}$.
\item ¿Cuántos grados de libertad tiene el numerador y cuántos el denominador?
\item Si en cambio se calcula $\dfrac{S_2^2}{S_1^2}$, ¿qué distribución tiene?
\end{enumerate}
\end{ejercicio}

\begin{solucion}
\textbf{(a)} Sabemos que $\dfrac{(n_1-1)S_1^2}{\sigma^2}\sim\chi^2_{n_1-1}=\chi^2_7$ y $\dfrac{(n_2-1)S_2^2}{\sigma^2}\sim\chi^2_{n_2-1}=\chi^2_{11}$, independientes (provienen de muestras independientes). Por definición de la $F$ de Snedecor,
\[
F = \frac{\left[\dfrac{(n_1-1)S_1^2}{\sigma^2}\right]\big/7}{\left[\dfrac{(n_2-1)S_2^2}{\sigma^2}\right]\big/11} = \frac{S_1^2}{S_2^2} \sim F_{7,11}.
\]

\textbf{(b)} Numerador: 7 grados de libertad; denominador: 11 grados de libertad.

\textbf{(c)} Por la propiedad de inversión de la $F$, $\dfrac{S_2^2}{S_1^2} = \dfrac{1}{F} \sim F_{11,7}$ (se invierten los grados de libertad).
\end{solucion}

\begin{ejercicio}
Una máquina produce piezas que se clasifican en 4 categorías de calidad: excelente (50\%), buena (30\%), regular (15\%) y defectuosa (5\%). Se inspeccionan 8 piezas producidas de forma independiente.
\begin{enumerate}[label=(\alph*)]
\item ¿Qué distribución tiene el vector $(X_1,X_2,X_3,X_4)$ con el número de piezas en cada categoría?
\item Calcular la probabilidad de obtener exactamente 4 excelentes, 2 buenas, 1 regular y 1 defectuosa.
\item Calcular $E[X_4]$, el número esperado de piezas defectuosas.
\end{enumerate}
\end{ejercicio}

\begin{solucion}
\textbf{(a)} $(X_1,X_2,X_3,X_4)$ tiene distribución \textbf{multinomial} con $n=8$ y probabilidades $p_1=0.50$, $p_2=0.30$, $p_3=0.15$, $p_4=0.05$.

\textbf{(b)}
\[
P(4,2,1,1) = \frac{8!}{4!\,2!\,1!\,1!}\,(0.50)^4(0.30)^2(0.15)^1(0.05)^1.
\]
Calculamos $\dfrac{8!}{4!\,2!\,1!\,1!} = \dfrac{40320}{24\cdot2\cdot1\cdot1} = 840$. Luego
\[
(0.50)^4=0.0625, \quad (0.30)^2=0.09, \quad 0.0625\times0.09\times0.15\times0.05 = 4.21875\times10^{-5}.
\]
\[
P(4,2,1,1) = 840 \times 4.21875\times10^{-5} \approx 0.03544.
\]

\textbf{(c)} Marginalmente $X_4\sim\text{Bin}(8, 0.05)$, entonces $E[X_4] = 8\times0.05 = 0.4$.
\end{solucion}

\begin{ejercicio}
Sean $X,Y$ con densidad conjunta uniforme en el cuadrado $(0,1)\times(0,1)$: $f_{X,Y}(x,y)=1$. Se definen $S=X+Y$ y $D=X-Y$ (cambio de variable a un nuevo sistema de funciones aleatorias).
\begin{enumerate}[label=(\alph*)]
\item Escribir la transformación inversa y calcular el jacobiano.
\item Hallar la densidad conjunta $f_{S,D}(s,d)$ y su soporte.
\end{enumerate}
\end{ejercicio}

\begin{solucion}
\textbf{(a)} De $s=x+y$, $d=x-y$ se despeja $x = \dfrac{s+d}{2}$, $y=\dfrac{s-d}{2}$. El jacobiano es
\[
J = \det\begin{pmatrix} \partial x/\partial s & \partial x/\partial d \\ \partial y/\partial s & \partial y/\partial d\end{pmatrix} = \det\begin{pmatrix} 1/2 & 1/2 \\ 1/2 & -1/2\end{pmatrix} = \left(\tfrac12\right)\left(-\tfrac12\right) - \left(\tfrac12\right)\left(\tfrac12\right) = -\tfrac14-\tfrac14 = -\tfrac12,
\]
luego $|J|=1/2$.

\textbf{(b)} Como $f_{X,Y}(x,y)=1$ en $(0,1)^2$,
\[
f_{S,D}(s,d) = f_{X,Y}\left(\tfrac{s+d}{2},\tfrac{s-d}{2}\right)\cdot|J| = 1\cdot\tfrac12 = \tfrac12,
\]
válida cuando $0<\tfrac{s+d}{2}<1$ y $0<\tfrac{s-d}{2}<1$, es decir, sobre el \textbf{cuadrado rotado} (rombo) de vértices $(0,0),(1,1),(2,0),(1,-1)$ en el plano $(s,d)$: $|d|<s$ para $0<s\le1$ y $|d|<2-s$ para $1<s<2$.
\end{solucion}

\end{document}
